% Sec. VIII Experiments (about 1.6 pages): spec sheet; E1-E4 test the
% instrument; E5 tests what a fitted channel misses.
\section{Experiments}
\label{sec:experiments}
Six tests check the instrument and what it adds (Table~\ref{tab:spec}).
E1 to E4 ask whether the testbed can be trusted; E5 asks whether a
synthetic channel with the same statistics reproduces the logged link; E6
asks whether the twin matches reality and whether what it trains
holds on real data. Each
pass criterion was stated before its test was run. Unless stated
otherwise, link numbers come from the 24 September session at site 1:
\texttt{survey\_1}, in which one agent owns the channel, and
\texttt{coop\_2}, in which both robots move and their active tests
contend. The pose reference is the 15 October mapping pass.

\begin{table*}[!t]
\caption{SCooP-iT Specification Sheet: Every Guarantee, Its Measured Bound, and the Check That Produced It}
\label{tab:spec}
\begin{center}
\renewcommand{\arraystretch}{1.2}
\setlength{\tabcolsep}{4pt}
\scriptsize
\begin{tabular}{@{}p{3.6cm} p{5.2cm} p{5.4cm} p{1.6cm}@{}}
\toprule
\textbf{Guarantee} & \textbf{Bound} & \textbf{Validated by} & \textbf{Experiment} \\
\midrule
Clock agreement between agents & measured offset $\leq$ \SI{0.99}{\milli\second} (\SI{1}{\milli\metre} at \SI{1}{\metre\per\second}); strict chrony bound $\leq$ \SI{12}{\milli\second}; reported RMS offsets \SIrange{10.9}{42.5}{\micro\second} per client & chrony offset logs, every run (coop\_2 so far) & E1 \\
Sniffer clock bias & \todo{to be measured} (worst case \SI{7.8}{\milli\second}) & probe-packet stamps on both ends & E1 \\
Spatial extrinsics & residual per sensor pair, Table~\ref{tab:calib} & held-out reflector and board positions & E1 \\
Agent pose in the site frame & \SI{1.06}{\centi\metre} surface agreement across passes (scale $2\times10^{-5}$); held-out boards \SI{10.0}{\milli\metre} / 0.88$^\circ$ (LiDAR robot), \SI{54.0}{\milli\metre} / 1.38$^\circ$ (RGB-D cart, coop2\_0828) and \SI{3.1}{\centi\metre} median (RGB-D cart, coop\_2); control distances \todo{} & second mapping pass, held-out boards, rangefinder control distances & E1, E2 \\
Link state & \SI{5}{\hertz} passive, \SI{100}{\hertz} RTT, \SI{4}{\second} goodput windows; outages labelled explicitly (3 outages, 12\,\% of frames on coop2\_0828) & monitor publishes through outages; map clearance predicts RSSI ($\rho = 0.60$, \SI{5.4}{\decibel}) & E2, E4 \\
Channel state (CSI) & \SIrange{104}{141}{\hertz}; valid below \SI{1.2}{\metre\per\second} relative speed & $\lambda/2$ sampling argument & --- \\
Contention as a knob & RTT rises several-fold idle to loaded (\SIrange{5}{10}{} to \SIrange{25}{33}{\milli\second}, 24 Sep; \SI{3.3}{} to \SIrange{60}{109}{\milli\second}, earlier pass); in-frame share of \SI{1.5}{\mega\byte} messages 86\,\% alone, \SIrange{59}{68}{\percent} with two robots & survey vs.\ cooperation mode; \todo{repeat across days} & E4, E5 \\
Replay fidelity & \todo{arrival-time error; frame-assignment agreement} & twin run: live payload against replayed trace & E3 \\
\bottomrule
\end{tabular}
\end{center}
\end{table*}


\subsection{E1 Precision: Errors Against Breaking Points}
\emph{Question.} Are the testbed's own errors small next to the errors
at which collaborative fusion breaks? \emph{Pass.} Every term at least
ten times below the breaking point of its kind.

Breaking points are derived for the objects the testbed scores, not
taken from vehicle studies. For a person, a
\SI{0.6}{}$\times$\SI{0.6}{\metre} box in bird's-eye view, a pose error
$d$ along one axis gives an IoU of $0.6(0.6-d)/[0.72-0.6(0.6-d)]$: 0.5 at
$d = \SI{0.20}{\metre}$ and 0.7 at $d = \SI{0.11}{\metre}$. Time errors
are converted to displacement at the measured 95th-percentile speed and
judged against the same points. Each term is plotted on one log axis
against both (Fig.~\ref{fig:headroom}).

\emph{Result.} At IoU 0.5 the measured clock offset (\SI{0.99}{\milli\second},
\SI{1}{\milli\metre} at \SI{1}{\metre\per\second}), the LiDAR robot's
held-out board residual (\SI{10}{\milli\metre}, $20\times$), and the
board detector (\SI{8}{\milli\metre}, $25\times$) pass, and the control
distances (\SI{2}{\centi\metre} RMS, $10\times$) meet the criterion
exactly. Three terms fall
short: the RGB-D cart (\SI{3.1}{\centi\metre} median, $6.5\times$), the
second-pass registration (\SI{3.0}{\centi\metre} RMS, $6.7\times$), and
the strict clock bound (\SI{12}{\milli\second}, $8.3\times$ against one
frame). At IoU 0.7 only the clock, the LiDAR robot, and the detector stay
near $10\times$, and the control distances fall to $5.5\times$; the cart's poses are then only about $3\times$ below the
breaking point, so perception scored from the cart's viewpoint at strict
IoU is limited by the reference. \todo{Fill $v_{95}$, the sniffer bias,
the LiDAR robot on \texttt{coop\_2}, and the radar residuals.}

\begin{figure}[t]
\centering
\fbox{\parbox[c][4.5cm][c]{0.95\columnwidth}{\centering
\todo{Fig.: E1 Headroom. One log axis in metres. Testbed error terms
(clock $\times$ speed, sniffer bias $\times$ speed, extrinsic residuals,
held-out board residual, second-pass agreement) against the fusion
breaking points (\SI{0.2}{\metre} pose error; latency as p95 speed
$\times$ delay against indoor object size \SIrange{0.5}{0.7}{\metre}).}}}
\caption{E1 Headroom: the testbed's own errors against the errors at
which collaborative fusion breaks.}
\label{fig:headroom}
\end{figure}


\subsection{E2 Consistency: The Map Predicts the Radio}
\emph{Question.} Are the poses and the link logs right if neither is
checked by hand? \emph{Pass.} Spearman $\rho > 0.3$ between Fresnel
clearance and the detrended signal, and a clear-minus-obstructed drop
above \SI{5}{\decibel}, on at least one link, with a 95\,\% block-bootstrap
interval excluding zero.

For every window, the first-Fresnel-zone clearance between the two ends
is computed through the map; a path is conventionally clear when
$0.6\,r_1$ of the first Fresnel radius $r_1 = \sqrt{\lambda d_1 d_2/d}$ is
free \cite{itup530}. The distance trend is removed with a fixed
path-loss exponent and an offset from the clear windows; fixing the
exponent keeps the wall loss out of the slope, since along a route
distance and obstruction move together. Neither the map nor the poses
used the signal, so agreement checks both. Intervals use a block
bootstrap over contiguous windows, since neighbouring windows are
correlated. Two confounds are controlled: the cart is pushed by a person
whose body can shadow its antenna, so windows are split by which side of
the cart faces the receiver; and the exponent is varied from 2 to 3.5.

\emph{Result.} On \texttt{coop\_2}, the channel from \texttt{mobile\_2}
to \texttt{mobile\_1}'s position, overheard by the sniffer, gives
$\rho = 0.60$ and a \SI{5.4}{\decibel} drop; the channel from the fixed
\texttt{infra\_1} gives $\rho = 0.23$ and \SI{4.7}{\decibel} and fails,
consistent with a mast whose path to the moving receiver changes less.
A self-test recovers a planted \SI{12}{\decibel} wall loss as
\SI{11.9}{\decibel}. \todo{Bootstrap intervals; exponent sweep; split by
cart side; pooled over \texttt{survey\_1} and \texttt{coop\_2}.}

\begin{figure}[t]
\centering
\fbox{\parbox[c][4cm][c]{0.95\columnwidth}{\centering
\todo{Fig.: E2 Map Meets Radio. Residual signal (dB, distance trend
removed with a fixed path-loss exponent) against Fresnel clearance, per
link and pooled; clear ($\geq 0.6$) vs.\ obstructed ($< 0.6$) windows.}}}
\caption{E2 Map Meets Radio: map-derived Fresnel clearance against the
measured signal it never used.}
\label{fig:clearance}
\end{figure}


\subsection{E3 Replay: Live Arrivals Against Replayed Ones}
\emph{Question.} Does replaying the logged link reproduce the arrivals a
method would see live? \emph{Pass.} At least 90\,\% agreement on the
receiver frame (\SI{100}{\milli\second}) in which each message lands, and
the logged channel closer than the fitted one.

On the \texttt{survey\_1} route, run A records the link as usual, and
run B, back to back, replaces the active tests with a payload node that
sends stamped messages at \SI{10}{\hertz} of \SI{1}{\kilo\byte},
\SI{100}{\kilo\byte}, and \SI{1.5}{\mega\byte} over the transport a
collaborative method uses (ROS~2, best effort), alone and with the other
robot sending; the same pair is repeated in cooperation mode. The
synchronized clocks give each message's one-way delay. B's messages are
replayed through A's trace and through the fitted and published
channels, and the arrival error, frame agreement, and drop agreement are
reported per size and mode (Fig.~\ref{fig:simgap}, left). If a size or
mode fails, the fluid model of Eq.~\eqref{eq:replay} is the limitation
and a per-packet trace \cite{nemfi} is the remedy. \emph{Result.}
\todo{Run the session.}

\subsection{E4 Repeatability: Same Setting, Same Link}
\emph{Question.} Does the same route and mode give the same link on
another day? \emph{Pass.} The day-to-day Wasserstein distance between
the RTT and goodput distributions is below half the distance between
survey and cooperation mode on the same route. \emph{Result.}
\todo{Mode distance from \texttt{survey\_1} and \texttt{coop\_2}; day
distance from the E3 session.}

\subsection{E5 Gap: What a Fitted Channel Misses}
\emph{Question.} Does a synthetic channel with the logged link's
statistics deliver messages as the logged link does? The fitted and
logged channels are compared, and the published channels are shown for
reference; there is no pass criterion, and either outcome is reported.
What the gap costs collaborative SLAM and perception, across methods
and sites, is measured on the SCooP dataset \cite{scoop_dataset}.

\emph{Setup and result.} Messages of \SI{1}{\kilo\byte}, \SI{100}{\kilo\byte},
and \SI{1.5}{\mega\byte} are sent through each channel at \SI{10}{\hertz}.
Under the logged link the predicted share of \SI{1.5}{\mega\byte} messages
that arrive within one frame is 86\,\% with one robot on the channel and
\SIrange{59}{68}{\percent} with two, while the loaded round trip,
\SIrange{25}{33}{\milli\second}, stays well inside the frame. The published
channels deliver 100\,\% by construction: V2X-ViT's constant
\SI{100}{\milli\second} sits on the inclusive frame boundary, and
Where2comm has no delay. These numbers are predictions of
Eq.~\eqref{eq:replay} until E3 validates it, and two robots sending raw
clouds at \SI{10}{\hertz} offer \SI{240}{\mega\bit\per\second}, close to the
measured goodput, so the share is also reported against offered load
divided by goodput, to separate a capacity limit from the link's
structure. \todo{Fitted column; share against normalized load.}

\begin{figure*}[t]
\centering
\fbox{\parbox[c][4.5cm][c]{0.95\textwidth}{\centering
\todo{Left (E3): replayed against live arrival per message, by size and
mode, for the logged and fitted channels. Right (E5a): in-frame share of
1\,KB, 100\,KB, and 1.5\,MB messages under the published, fitted, and
logged channels, survey and cooperation mode, against offered load over
goodput.}}}
\caption{Left, E3: replayed against live message arrivals. Right, E5a:
the logged link against a synthetic channel with the same statistics and
against the published channels of prior benchmarks.}
\label{fig:simgap}
\end{figure*}


\subsection{E6 Twin: The Twin Against Held-Out Passes}
\emph{Question.} Does the twin reproduce held-out real passes?
\emph{Pass.} Synthetic
scans at the recorded poses within \SI{3}{\centi\metre} median of the real
scans; fitted dropout rates and burst lengths within the spread between
real sessions; the place-indexed link predicting the held-out pass's
in-frame share within 10 points and closer than the fitted channel.

\emph{Setup.} The twin drives the recorded paths of a held-out
pass; its scans are registered to the map and compared with the real
scans at the same poses, its camera views are compared with the
recorded images at the same poses, and the dropout and link models fitted on the
build passes are compared with that pass. Whether models trained in
the twin hold on real recordings needs labels and is reported with the
SCooP dataset \cite{scoop_dataset}. \emph{Result.} \todo{Run the twin
on a held-out pass.}

